% ============================================================
\chapter{M\'ethodes d'Ensemble --- Bagging, Boosting, Random Forests}
\label{ch:ensembles}
\index{m\'ethodes d'ensemble}
% ============================================================

\section{Pourquoi les m\'ethodes d'ensemble~?}
\label{sec:ensembles-pourquoi}

L'id\'ee fondamentale des m\'ethodes d'ensemble est de combiner plusieurs
mod\`eles <<~faibles~>> pour obtenir un mod\`ele <<~fort~>>. Nous montrons
ici formellement pourquoi cette strat\'egie r\'eduit l'erreur.

\begin{theoreme}[R\'eduction de variance par moyennage]
\label{thm:variance-reduction}
Soient $f_1, \ldots, f_B$ des mod\`eles ind\'ependants et identiquement
distribu\'es, chacun de variance $\sigma^2$ et de biais $\mu$. Le mod\`ele
moyenn\'e $\bar{f} = \frac{1}{B}\sum_{b=1}^B f_b$ satisfait~:
\begin{equation}
  \Bias[\bar{f}] = \mu, \qquad \Var[\bar{f}] = \frac{\sigma^2}{B}
\end{equation}
Le biais est inchang\'e mais la variance est divis\'ee par $B$.
\end{theoreme}

\begin{proof}
Par lin\'earit\'e de l'esp\'erance~:
$\E[\bar{f}] = \frac{1}{B}\sum_b \E[f_b] = \mu$. Par ind\'ependance~:
$\Var[\bar{f}] = \frac{1}{B^2}\sum_b \Var[f_b] = \frac{\sigma^2}{B}$.
\end{proof}

\begin{remarque}
En pratique, les mod\`eles ne sont pas parfaitement ind\'ependants car ils
sont entra\^in\'es sur des \'echantillons li\'es. Si la corr\'elation entre
paires est $\rho$, alors~:
\begin{equation}
  \Var[\bar{f}] = \rho\sigma^2 + \frac{1 - \rho}{B}\sigma^2
  \label{eq:var-correlated}
\end{equation}
R\'eduire $\rho$ est la cl\'e pour am\'eliorer les ensembles.
\end{remarque}

\begin{center}
\begin{tikzpicture}[
  box/.style={rectangle, draw, rounded corners, minimum width=2cm,
    minimum height=0.7cm, align=center, font=\small},
  >=Stealth
]
  % Bagging
  \node[box, fill=blue!10] (bag) at (0, 0) {\textbf{Bagging}};
  \node[box, fill=blue!5, below left=0.8cm and 0.5cm of bag] (b1)
    {Mod\`ele 1};
  \node[box, fill=blue!5, below=0.8cm of bag] (b2) {Mod\`ele 2};
  \node[box, fill=blue!5, below right=0.8cm and 0.5cm of bag] (b3)
    {Mod\`ele $B$};
  \node[box, fill=green!15, below=0.8cm of b2] (bavg) {Moyenne / Vote};
  \draw[->] (bag) -- (b1);
  \draw[->] (bag) -- (b2);
  \draw[->] (bag) -- (b3);
  \draw[->] (b1) -- (bavg);
  \draw[->] (b2) -- (bavg);
  \draw[->] (b3) -- (bavg);
  \node[font=\scriptsize, right=0.1cm of b2] {$\cdots$};
  \node[font=\scriptsize, above=0.1cm of bag] {Parall\`ele};

  % Boosting
  \node[box, fill=orange!10] (boost) at (8, 0) {\textbf{Boosting}};
  \node[box, fill=orange!5, below=0.8cm of boost] (t1) {Mod\`ele 1};
  \node[box, fill=orange!5, below=0.8cm of t1] (t2) {Mod\`ele 2};
  \node[box, fill=orange!5, below=0.8cm of t2] (t3) {Mod\`ele $T$};
  \node[box, fill=green!15, right=1.5cm of t2] (bsum) {Somme pond\'er\'ee};
  \draw[->] (boost) -- (t1);
  \draw[->] (t1) -- node[right, font=\scriptsize] {r\'esidus} (t2);
  \draw[->] (t2) -- node[right, font=\scriptsize] {r\'esidus} (t3);
  \node[font=\scriptsize, left=0.2cm of t2] {$\vdots$};
  \draw[->] (t1.east) -- ++(0.5,0) |- (bsum);
  \draw[->] (t2) -- (bsum);
  \draw[->] (t3.east) -- ++(0.5,0) |- (bsum);
  \node[font=\scriptsize, above=0.1cm of boost] {S\'equentiel};
\end{tikzpicture}
\end{center}

\section{Bagging (Bootstrap Aggregating)}
\label{sec:bagging}
\index{bagging}

\begin{definition}[Bagging]
Le \textbf{bagging}~\cite{hastie2009} consiste \`a~:
\begin{enumerate}
  \item G\'en\'erer $B$ \'echantillons bootstrap $\mathcal{D}_1^*, \ldots,
    \mathcal{D}_B^*$ par tirage avec remise de taille $n$.
  \item Entra\^iner un mod\`ele $f_b$ sur chaque $\mathcal{D}_b^*$.
  \item Agr\'eger les pr\'edictions~:
    \begin{itemize}
      \item R\'egression~:
        $\hat{f}(\bm{x}) = \frac{1}{B}\sum_{b=1}^B f_b(\bm{x})$
      \item Classification~:
        $\hat{f}(\bm{x}) = \argmax_k \sum_{b=1}^B \ind[f_b(\bm{x}) = k]$
        (vote majoritaire)
    \end{itemize}
\end{enumerate}
\end{definition}

\begin{proposition}[Proportion d'exemples hors-sac]
\index{bootstrap}
Un \'echantillon bootstrap de taille $n$ contient en moyenne environ
$1 - (1 - 1/n)^n \approx 1 - e^{-1} \approx 63.2\,\%$ des exemples
originaux. Les $\sim 36.8\,\%$ restants sont dits <<~hors-sac~>>
(\emph{out-of-bag}).
\end{proposition}

\begin{proof}
La probabilit\'e qu'un exemple $i$ ne soit pas s\'electionn\'e dans un tirage
est $1 - 1/n$. Sur $n$ tirages ind\'ependants~:
$\Pr(i \notin \mathcal{D}_b^*) = (1 - 1/n)^n \xrightarrow[n \to \infty]{}
e^{-1}$.
\end{proof}

\subsection{D\'erivation formelle de la r\'eduction de variance}

\begin{theoreme}[Variance du bagging]
Si les mod\`eles $f_b$ ont une variance $\sigma^2$ et une corr\'elation par
paires $\rho$, alors la pr\'ediction bag\'ee a pour variance~:
\begin{equation}
  \Var[\hat{f}_{\text{bag}}] = \rho\sigma^2
  + \frac{(1-\rho)\sigma^2}{B}
\end{equation}
Pour $B \to \infty$, $\Var[\hat{f}_{\text{bag}}] \to \rho\sigma^2$.
\end{theoreme}

\begin{proof}
On a $\hat{f}_{\text{bag}} = \frac{1}{B}\sum_b f_b$. Donc~:
\begin{align}
  \Var[\hat{f}_{\text{bag}}]
  &= \frac{1}{B^2}\left[\sum_b \Var[f_b]
    + \sum_{b \neq b'} \Cov(f_b, f_{b'})\right] \\
  &= \frac{1}{B^2}\left[B\sigma^2 + B(B-1)\rho\sigma^2\right] \\
  &= \frac{\sigma^2}{B} + \frac{(B-1)\rho\sigma^2}{B}
  = \rho\sigma^2 + \frac{(1-\rho)\sigma^2}{B}
\end{align}
\end{proof}

\section{For\^ets al\'eatoires}
\label{sec:random-forests}
\index{for\^ets al\'eatoires}

\begin{definition}[For\^et al\'eatoire]
Une \textbf{for\^et al\'eatoire}~\cite{hastie2009} est un bagging d'arbres de
d\'ecision avec une randomisation suppl\'ementaire~: \`a chaque s\'eparation,
seul un sous-ensemble al\'eatoire de $m$ variables (sur $d$) est
consid\'er\'e.
\end{definition}

\begin{formules}[title=\textbf{Choix de $m$ (nombre de variables candidates)}]
\begin{center}
\begin{tabular}{lc}
\toprule
\textbf{T\^ache} & \textbf{Valeur recommand\'ee} \\
\midrule
Classification & $m = \lfloor\sqrt{d}\rfloor$ \\
R\'egression & $m = \lfloor d/3 \rfloor$ \\
\bottomrule
\end{tabular}
\end{center}
R\'eduire $m$ diminue la corr\'elation $\rho$ entre arbres, ce qui r\'eduit
la variance de l'ensemble (cf. \'equation~\eqref{eq:var-correlated}).
\end{formules}

\begin{algorithme}[title=\textbf{Algorithme Random Forest}]
\textbf{Entr\'ee}~: $\mathcal{D} = \{(\bm{x}_i, y_i)\}_{i=1}^n$, nombre
d'arbres $B$, nombre de variables $m$.

\textbf{Pour} $b = 1, \ldots, B$~:
\begin{enumerate}
  \item Tirer un \'echantillon bootstrap $\mathcal{D}_b^*$ de taille $n$.
  \item Construire un arbre $T_b$ sur $\mathcal{D}_b^*$~:
    \begin{itemize}
      \item \`A chaque n\oe{}ud, s\'electionner al\'eatoirement $m$ variables.
      \item Trouver la meilleure s\'eparation parmi ces $m$ variables.
      \item D\'evelopper l'arbre compl\`etement (pas d'\'elagage).
    \end{itemize}
\end{enumerate}
\textbf{Pr\'ediction}~: agr\'eger par vote (classification) ou moyenne
(r\'egression).
\end{algorithme}

\subsection{Erreur out-of-bag (OOB)}
\index{erreur OOB}

\begin{definition}[Erreur OOB]
Pour chaque observation $(\bm{x}_i, y_i)$, l'erreur \textbf{out-of-bag} est
calcul\'ee en agr\'egeant uniquement les arbres pour lesquels $i$ n'\'etait
pas dans l'\'echantillon bootstrap~:
\begin{equation}
  \hat{y}_i^{\text{OOB}} = \text{agr\'eger}\{T_b(\bm{x}_i) :
    i \notin \mathcal{D}_b^*\}
\end{equation}
L'erreur OOB globale est une estimation non biais\'ee de l'erreur de
g\'en\'eralisation, sans n\'ecessiter d'ensemble de validation.
\end{definition}

\begin{bonnepratique}[title=\textbf{Avantages des for\^ets al\'eatoires}]
\begin{itemize}
  \item Peu d'hyperparam\`etres \`a r\'egler (principalement $B$ et $m$).
  \item Robustes au surapprentissage~: augmenter $B$ ne d\'egrade jamais
    la performance (la variance diminue ou stagne).
  \item Fournissent une estimation de l'erreur (OOB) gratuitement.
  \item Calculent l'importance des variables.
  \item Parall\'elisables efficacement.
\end{itemize}
\end{bonnepratique}

\section{Boosting}
\label{sec:boosting}
\index{boosting}

Contrairement au bagging (r\'eduction de variance), le boosting vise \`a
r\'eduire le \textbf{biais} en combinant s\'equentiellement des mod\`eles
faibles.

\subsection{AdaBoost}
\index{AdaBoost}

\begin{algorithme}[title=\textbf{Algorithme AdaBoost}]
\textbf{Entr\'ee}~: $\{(\bm{x}_i, y_i)\}_{i=1}^n$ avec $y_i \in \{-1, +1\}$,
nombre d'it\'erations $T$.

\textbf{Initialisation}~: $w_i^{(1)} = 1/n$ pour tout $i$.

\textbf{Pour} $t = 1, \ldots, T$~:
\begin{enumerate}
  \item Entra\^iner un classifieur faible $h_t$ en minimisant l'erreur
    pond\'er\'ee~:
    \begin{equation}
      \varepsilon_t = \sum_{i=1}^n w_i^{(t)} \ind[h_t(\bm{x}_i) \neq y_i]
    \end{equation}
  \item Calculer le poids du classifieur~:
    \begin{equation}
      \alpha_t = \frac{1}{2}\ln\!\left(\frac{1 - \varepsilon_t}{\varepsilon_t}
      \right)
      \label{eq:adaboost-alpha}
    \end{equation}
  \item Mettre \`a jour les poids des exemples~:
    \begin{equation}
      w_i^{(t+1)} = \frac{w_i^{(t)} \exp(-\alpha_t y_i h_t(\bm{x}_i))}{Z_t}
      \label{eq:adaboost-weights}
    \end{equation}
    o\`u $Z_t$ est le facteur de normalisation.
\end{enumerate}
\textbf{Pr\'ediction}~:
$H(\bm{x}) = \sign\!\left(\sum_{t=1}^T \alpha_t h_t(\bm{x})\right)$
\end{algorithme}

\begin{theoreme}[D\'erivation de la r\`egle de mise \`a jour]
\index{AdaBoost!d\'erivation}
AdaBoost minimise la perte exponentielle~:
\begin{equation}
  L = \sum_{i=1}^n \exp\!\left(-y_i \sum_{t=1}^T \alpha_t h_t(\bm{x}_i)\right)
\end{equation}
\`A l'\'etape $t$, en fixant les classifieurs pr\'ec\'edents et en optimisant
$\alpha_t$~:
\begin{align}
  \frac{\partial L}{\partial \alpha_t}
  &= -\sum_{i=1}^n y_i h_t(\bm{x}_i) w_i^{(t)} e^{-\alpha_t y_i h_t(\bm{x}_i)}
  = 0
\end{align}
En s\'eparant les exemples correctement class\'es ($y_i h_t = +1$) et
incorrectement class\'es ($y_i h_t = -1$)~:
\begin{equation}
  e^{-\alpha_t}(1 - \varepsilon_t) - e^{\alpha_t}\varepsilon_t = 0
  \implies \alpha_t = \frac{1}{2}\ln\!\left(\frac{1 - \varepsilon_t}
  {\varepsilon_t}\right)
\end{equation}
\end{theoreme}

\begin{proposition}[Borne sur l'erreur d'entra\^inement]
L'erreur d'entra\^inement d'AdaBoost est born\'ee par~:
\begin{equation}
  \frac{1}{n}\sum_{i=1}^n \ind[H(\bm{x}_i) \neq y_i]
  \leq \prod_{t=1}^T 2\sqrt{\varepsilon_t(1 - \varepsilon_t)}
  = \prod_{t=1}^T \sqrt{1 - 4\gamma_t^2}
\end{equation}
o\`u $\gamma_t = \frac{1}{2} - \varepsilon_t$ est l'<<~avantage~>> du
classifieur faible. Si $\gamma_t \geq \gamma > 0$, l'erreur
d\'ecro\^it exponentiellement en $T$.
\end{proposition}

\begin{attention}[title=\textbf{AdaBoost et le bruit}]
AdaBoost augmente les poids des exemples mal class\'es. En pr\'esence de bruit
(\'etiquettes erron\'ees), il sur-pond\`ere les outliers, ce qui peut
d\'egrader la g\'en\'eralisation. Le Gradient Boosting avec r\'egularisation
est g\'en\'eralement plus robuste.
\end{attention}

\subsection{Gradient Boosting}
\index{gradient boosting}

\begin{definition}[Gradient Boosting]
Le \textbf{gradient boosting} g\'en\'eralise le boosting \`a toute fonction
de perte diff\'erentiable $\ell$. Le mod\`ele est construit de fa\c{c}on
additive~:
\begin{equation}
  F_t(\bm{x}) = F_{t-1}(\bm{x}) + \eta \cdot h_t(\bm{x})
\end{equation}
o\`u $h_t$ est ajust\'e sur les \textbf{pseudo-r\'esidus}
(gradient n\'egatif)~:
\begin{equation}
  r_i^{(t)} = -\frac{\partial \ell(y_i, F_{t-1}(\bm{x}_i))}
    {\partial F_{t-1}(\bm{x}_i)}
  \label{eq:pseudo-residuals}
\end{equation}
et $\eta \in (0, 1]$ est le taux d'apprentissage (\emph{learning rate}).
\end{definition}

\begin{algorithme}[title=\textbf{Algorithme Gradient Boosting}]
\textbf{Entr\'ee}~: $\{(\bm{x}_i, y_i)\}_{i=1}^n$, perte $\ell$, nombre
d'it\'erations $T$, taux $\eta$.

\begin{enumerate}
  \item Initialiser $F_0(\bm{x}) = \argmin_c \sum_{i=1}^n \ell(y_i, c)$.
  \item \textbf{Pour} $t = 1, \ldots, T$~:
    \begin{enumerate}
      \item Calculer les pseudo-r\'esidus $r_i^{(t)}$
        (\'equation~\eqref{eq:pseudo-residuals}).
      \item Ajuster un arbre de r\'egression $h_t$ sur
        $\{(\bm{x}_i, r_i^{(t)})\}$.
      \item Mettre \`a jour~: $F_t = F_{t-1} + \eta \cdot h_t$.
    \end{enumerate}
\end{enumerate}
\textbf{Pr\'ediction}~: $\hat{y} = F_T(\bm{x})$.
\end{algorithme}

\begin{intuition}[title=\textbf{Descente de gradient dans l'espace fonctionnel}]
Le gradient boosting effectue une descente de gradient, non pas dans l'espace
des param\`etres, mais dans l'\textbf{espace des fonctions}. \`A chaque
it\'eration, $h_t$ approxime la direction de plus forte descente de la perte.
Le taux $\eta$ joue le r\^ole du pas de gradient et r\'egularise le mod\`ele~:
un $\eta$ petit n\'ecessite plus d'arbres mais g\'en\'eralise mieux.
\end{intuition}

\subsection{XGBoost et LightGBM}
\index{XGBoost}\index{LightGBM}

Les impl\'ementations modernes ajoutent plusieurs am\'eliorations~:

\begin{formules}[title=\textbf{Am\'eliorations de XGBoost/LightGBM}]
\begin{center}
\begin{tabular}{lp{8cm}}
\toprule
\textbf{Technique} & \textbf{Description} \\
\midrule
R\'egularisation L1/L2 & P\'enalisation des poids des feuilles~:
  $\Omega(h) = \gamma|T| + \frac{\lambda}{2}\sum_j w_j^2$ \\
Approximation de Newton & Utilisation du hessien (d\'eriv\'ee seconde) pour
  des s\'eparations plus pr\'ecises \\
Sous-\'echantillonnage & Colonnes et lignes \'echantillonn\'ees \`a chaque
  arbre (r\'eduit $\rho$) \\
Croissance par feuille & LightGBM fait cro\^itre l'arbre par feuille
  (\emph{leaf-wise}) au lieu de par niveau (\emph{level-wise}) \\
Histogrammes & Discr\'etisation des features pour acc\'el\'erer la recherche
  de seuils \\
\bottomrule
\end{tabular}
\end{center}
\end{formules}

\section{Comparaison des m\'ethodes}
\label{sec:comparaison}

\begin{formules}[title=\textbf{Tableau comparatif}]
\begin{center}
\begin{tabular}{lccc}
\toprule
& \textbf{Bagging/RF} & \textbf{AdaBoost} & \textbf{Gradient Boosting} \\
\midrule
R\'eduit & Variance & Biais & Biais \\
Parall\`ele & Oui & Non & Non \\
Robuste au bruit & Oui & Non & Mod\'er\'e \\
Risque surapprentissage & Faible & Mod\'er\'e & \'Elev\'e \\
Hyperparam\`etres & Peu & Peu & Beaucoup \\
\bottomrule
\end{tabular}
\end{center}
\end{formules}

\section{Impl\'ementations avec scikit-learn}
\label{sec:ensembles-sklearn}

\begin{codeblock}[title=\textbf{Random Forest}]
\begin{minted}{python}
import numpy as np
from sklearn.datasets import load_wine
from sklearn.ensemble import RandomForestClassifier
from sklearn.model_selection import train_test_split, cross_val_score

X, y = load_wine(return_X_y=True)
X_train, X_test, y_train, y_test = train_test_split(
    X, y, test_size=0.3, random_state=42, stratify=y
)

rf = RandomForestClassifier(
    n_estimators=200, max_features="sqrt", oob_score=True,
    random_state=42, n_jobs=-1
)
rf.fit(X_train, y_train)

print("=== Random Forest ===")
print(f"Accuracy train : {rf.score(X_train, y_train):.4f}")
print(f"Accuracy test  : {rf.score(X_test, y_test):.4f}")
print(f"Erreur OOB     : {1 - rf.oob_score_:.4f}")
print(f"CV accuracy    : {cross_val_score(rf, X, y, cv=5).mean():.4f}")

# Importance des variables
importances = rf.feature_importances_
feature_names = load_wine().feature_names
indices = np.argsort(importances)[::-1]
print("\nTop 5 variables :")
for i in range(5):
    print(f"  {feature_names[indices[i]]:25s} : {importances[indices[i]]:.4f}")
\end{minted}
\end{codeblock}

\begin{codeoutput}
\begin{verbatim}
=== Random Forest ===
Accuracy train : 1.0000
Accuracy test  : 0.9815
Erreur OOB     : 0.0323
CV accuracy    : 0.9775

Top 5 variables :
  color_intensity           : 0.1542
  flavanoids                : 0.1487
  proline                   : 0.1325
  alcohol                   : 0.1198
  od280/od315_of_diluted_wines : 0.0843
\end{verbatim}
\end{codeoutput}

\begin{codeblock}[title=\textbf{Gradient Boosting et comparaison}]
\begin{minted}{python}
from sklearn.ensemble import (
    GradientBoostingClassifier, AdaBoostClassifier
)
from sklearn.tree import DecisionTreeClassifier
import matplotlib.pyplot as plt

# AdaBoost
ada = AdaBoostClassifier(
    estimator=DecisionTreeClassifier(max_depth=1),
    n_estimators=200, learning_rate=0.5, random_state=42
)
ada.fit(X_train, y_train)

# Gradient Boosting
gb = GradientBoostingClassifier(
    n_estimators=200, max_depth=3, learning_rate=0.1,
    subsample=0.8, random_state=42
)
gb.fit(X_train, y_train)

print("=== Comparaison ===")
models = {"Random Forest": rf, "AdaBoost": ada, "Gradient Boosting": gb}
for name, model in models.items():
    acc_train = model.score(X_train, y_train)
    acc_test = model.score(X_test, y_test)
    cv = cross_val_score(model, X, y, cv=5).mean()
    print(f"{name:20s} | train={acc_train:.4f} | "
          f"test={acc_test:.4f} | CV={cv:.4f}")
\end{minted}
\end{codeblock}

\begin{codeoutput}
\begin{verbatim}
=== Comparaison ===
Random Forest        | train=1.0000 | test=0.9815 | CV=0.9775
AdaBoost             | train=1.0000 | test=0.9444 | CV=0.9497
Gradient Boosting    | train=1.0000 | test=0.9815 | CV=0.9719
\end{verbatim}
\end{codeoutput}

\begin{codeblock}[title=\textbf{Courbes d'apprentissage et effet du nombre d'arbres}]
\begin{minted}{python}
# Erreur en fonction du nombre d'arbres (staged_predict)
n_estimators_range = range(1, 201)

fig, axes = plt.subplots(1, 2, figsize=(12, 5))

# AdaBoost
ada_train = [1 - acc for acc in ada.staged_score(X_train, y_train)]
ada_test = [1 - acc for acc in ada.staged_score(X_test, y_test)]
axes[0].plot(n_estimators_range, ada_train, label="Train")
axes[0].plot(n_estimators_range, ada_test, label="Test")
axes[0].set_title("AdaBoost")
axes[0].set_xlabel("Nombre d'arbres")
axes[0].set_ylabel("Erreur")
axes[0].legend()

# Gradient Boosting
gb_train = [1 - acc for acc in gb.staged_score(X_train, y_train)]
gb_test = [1 - acc for acc in gb.staged_score(X_test, y_test)]
axes[1].plot(n_estimators_range, gb_train, label="Train")
axes[1].plot(n_estimators_range, gb_test, label="Test")
axes[1].set_title("Gradient Boosting")
axes[1].set_xlabel("Nombre d'arbres")
axes[1].set_ylabel("Erreur")
axes[1].legend()

plt.tight_layout()
plt.savefig("ensemble_learning_curves.pdf")
\end{minted}
\end{codeblock}

\begin{codeblock}[title=\textbf{XGBoost (si disponible)}]
\begin{minted}{python}
try:
    from xgboost import XGBClassifier

    xgb = XGBClassifier(
        n_estimators=200, max_depth=4, learning_rate=0.1,
        subsample=0.8, colsample_bytree=0.8,
        reg_lambda=1.0, reg_alpha=0.1,
        eval_metric="mlogloss", random_state=42
    )
    xgb.fit(X_train, y_train)
    print(f"XGBoost test accuracy : {xgb.score(X_test, y_test):.4f}")
    print(f"XGBoost CV accuracy   : "
          f"{cross_val_score(xgb, X, y, cv=5).mean():.4f}")
except ImportError:
    print("xgboost non installe. Installer avec: pip install xgboost")
\end{minted}
\end{codeblock}

\begin{codeoutput}
\begin{verbatim}
XGBoost test accuracy : 0.9815
XGBoost CV accuracy   : 0.9775
\end{verbatim}
\end{codeoutput}

\section{Exercices}
\label{sec:exercises-ch7}

\begin{exercice}[Bootstrap et erreur OOB --- Niveau 1]
\begin{enumerate}
  \item Montrer que la probabilit\'e qu'un exemple donn\'e soit absent d'un
    \'echantillon bootstrap de taille $n$ converge vers $e^{-1}$ quand
    $n \to \infty$.
  \item Pour $n = 100$, calculer la valeur exacte
    $(1 - 1/100)^{100}$ et la comparer \`a $e^{-1}$.
  \item Entra\^iner une for\^et al\'eatoire sur le dataset Iris avec
    \texttt{oob\_score=True}. Comparer l'erreur OOB avec l'erreur de
    validation crois\'ee 5-fold.
\end{enumerate}
\end{exercice}

\begin{exercice}[AdaBoost analytique --- Niveau 2]
Soit un jeu de donn\'ees de 6 points \'equir\'epartis entre les classes
$+1$ et $-1$. \`A la premi\`ere it\'eration, le classifieur faible $h_1$
classe correctement 4 points sur 6.
\begin{enumerate}
  \item Calculer $\varepsilon_1$, $\alpha_1$ et les nouveaux poids $w_i^{(2)}$.
  \item Si le deuxi\`eme classifieur faible $h_2$ classe correctement les 2
    points mal class\'es par $h_1$ mais en rate 1 autre, calculer
    $\varepsilon_2$ et $\alpha_2$.
  \item D\'emontrer compl\`etement que la formule~\eqref{eq:adaboost-alpha}
    minimise la perte exponentielle \`a chaque \'etape.
  \item V\'erifier que la borne sur l'erreur d'entra\^inement est
    satisfaite apr\`es 2 it\'erations.
\end{enumerate}
\end{exercice}

\begin{exercice}[Comparaison compl\`ete --- Niveau 3]
\begin{enumerate}
  \item Sur le dataset \texttt{breast\_cancer} de scikit-learn, comparer
    syst\'ematiquement~: arbre seul, bagging (50 arbres), for\^et al\'eatoire
    (50, 200, 500 arbres), AdaBoost (100 stumps), Gradient Boosting
    (100 arbres, $\eta = 0.1$). Reporter accuracy, F1-score et temps
    d'entra\^inement (validation crois\'ee 10-fold).
  \item \'Etudier l'effet du nombre de variables candidates $m$ dans une
    for\^et al\'eatoire. Tracer l'erreur OOB en fonction de $m$ pour
    $m \in \{1, 2, \ldots, d\}$.
  \item Impl\'ementer l'algorithme AdaBoost \`a partir de z\'ero (sans
    scikit-learn) avec des \emph{decision stumps} comme classifieurs faibles.
    Comparer avec \texttt{AdaBoostClassifier}.
  \item D\'emontrer formellement que le bagging ne peut pas augmenter le biais
    par rapport au mod\`ele de base (pour la perte quadratique), en utilisant
    l'in\'egalit\'e de Jensen.
\end{enumerate}
\end{exercice}
